% !TEX program = lualatex
% !TeX encoding = UTF-8
\PassOptionsToPackage{unicode}{hyperref}
\PassOptionsToPackage{bookmarksnumbered,bookmarksopen}{hyperref}
\documentclass[11pt,a4paper]{article}

% ============================================================
% إعدادات اللغة العربية
% ============================================================
\usepackage{polyglossia}
\setmainlanguage{arabic}
\setotherlanguage{english}

% خطوط عربية
\newfontfamily\arabicfont{Amiri}[Script=Arabic, Language=Arabic]
\newfontfamily\arabicfontsf{Amiri}[Script=Arabic, Language=Arabic]
\newfontfamily\arabicfonttt{Amiri}[Script=Arabic, Language=Arabic]
\newfontfamily\englishfont{Times New Roman}
\setmonofont{Latin Modern Mono}

% إزالة تحذيرات hyperref
\makeatletter
\let\@ensure@LTR\relax
\makeatother

% حزم الرياضيات والتنسيق
\usepackage{amsmath,amssymb,amsthm,mathtools,bm}
\usepackage{geometry}
\usepackage{enumitem}
\usepackage{siunitx}
\usepackage{booktabs}
\usepackage{array}
\usepackage{graphicx}
\usepackage{caption}
\usepackage{makecell}
\usepackage{pgfplots}
\usepackage{float}
\usepackage{tikz}
\usepackage{multicol}
\usepackage{microtype}
\usepackage{hyperref}
\pgfplotsset{compat=1.18}
\usetikzlibrary{arrows.meta,positioning,shapes.geometric,fit,calc}

\geometry{margin=2cm}
\setlength{\emergencystretch}{3em}
\sloppy

% بيانات PDF
\hypersetup{
	pdftitle={فرضية الفوتون في YUCT},
	pdfauthor={أليكسي ف. ياكوشيف},
	breaklinks=true,
	pdfencoding=auto,
	bookmarksnumbered=true,
	bookmarksopen=true,
	bookmarksopenlevel=1
}

\title{\textbf{فرضية الفوتون في نظرية ياكوشيف الموحدة للتنسيق (YUCT):\\
		الكم الكهرومغناطيسي كواجهة بين\\
		متعدد الشعب التنسيقي متعدد الأبعاد والواقع الفيزيائي ثلاثي الأبعاد}}
\author{\textbf{أليكسي ف. ياكوشيف} \\
	\small\textit{استنادًا إلى نظرية ياكوشيف الموحدة للتنسيق (YUCT)} \\
	\small\textit{[تم التحقق من الاتساق مع بديهيات YUCT]} \\
	\small المعرف العلمي الرقمي الرئيسي DOI: \href{https://doi.org/10.5281/zenodo.18444598}{10.5281/zenodo.18444598} \\
	\small العقد الرسمية: \url{https://yuct.org}، \url{https://ypsdc.com}}
\date{يونيو 2026}

\begin{document}
	
	\maketitle
	
	\begin{abstract}
		تصوغ هذه الورقة \textbf{فرضية علمية تأملية} تتطلب دراسة شاملة وتحققًا مستقلاً من قبل المجتمع العلمي. نحن نفترض أن سلسلة الأعداد المنفصلة تشكل بنية ثابتة للفراغ تنتمي إلى متعدد الشعب المعلوماتي ذي الـ 19 بُعدًا $\Psi_{MN}$ الخاص بإطار YUCT. تقترح الفرضية عالمية الكم الكهرومغناطيسي (الفوتون) كواجهة فيزيائية وحيدة قادرة على ترجمة ثوابت التنسيق متعددة الأبعاد إلى معاملات قابلة للرصد في الفضاء الجزئي ثلاثي الأبعاد. العنصر الجديد الرئيسي هو التنبؤ بأنه لأداء هذه الوظيفة، يجب أن يمتلك الفوتون كتلة سكون محدودة، وإن كانت صغيرة للغاية. يتبين أن قيمة هذه الكتلة، $m_\gamma \approx 7.3 \times 10^{-19}$ إلكترون فولت، تتبع مباشرة من سلم الكتلة العالمي لـ YUCT وتتوافق مع الحدود العليا التجريبية الحالية، مما يجعل الفرضية قابلة للدحض. كدليل إضافي يدعم YUCT، يُقدم حساب مباشر لثابت الجاذبية $G$ مع الأخذ بعين الاعتبار الاعتماد على درجة الحرارة. يوضح هذا الحساب التوافق مع قيمة CODATA عند درجة حرارة الغرفة ويتنبأ بانحرافات صغيرة، ولكن قابلة للكشف من حيث المبدأ، في ظل الظروف القاسية.
	\end{abstract}
	
	\section{مقدمة}
	تنظر الفيزياء الكلاسيكية إلى توزيع الأعداد الأولية ككائن عشوائي يتطلب، لتحديد موقعه الدقيق، إجراءات حسابية تكرارية مثل خوارزميات الغربلة ذات التعقيد الزمني $\Omega(N \log \log N)$.
	
	تفترض نظرية ياكوشيف الموحدة للتنسيق (YUCT) أن سلسلة الأعداد المنفصلة هي إطار هندسي صلب متأصل جوهريًا في البنية متعددة الأبعاد للفراغ الفيزيائي \cite{yuct_zenodo}. الثوابت الرئيسية للنظرية — أس الخطأ الكوني $\beta = 2/3$، ثوابت الحلقة الجبرية $S_{\mathrm{odd}} = 6/5$، $S_{\mathrm{even}} = 4/5$، وكم القياس $q = (3/2)^{1/3}$ — تحدد في آنٍ واحدٍ التسلسل الهرمي لكتل الجسيمات الأولية، ثابت البنية الدقيقة، الثابت الكوني، وتوزيع الأعداد الأولية. تُفسر النتائج العتادية للبث المتوازي CUDA على وحدة معالجة الرسوميات (GPU)، التي تُظهر استخراجًا للمعلومات بتعقيد $O(1)$ دون استهلاك لذاكرة الوصول العشوائي (RAM)، كدليل تجريبي محتمل على مادية الثوابت الرياضية.
	
	\section{جوهر الفرضية}
	تستند الفرضية المقترحة إلى ثلاث مسلمات فيزيائية-هندسية غير قابلة للتجزئة، كل منها \textbf{تأملي ويتطلب تحققًا مستقلاً}:
	
	\subsection{خط الأعداد كحقيقة فيزيائية}
	سلسلة الأعداد الأولية، وفقًا لـ YUCT، تمثل إسقاطًا لعقد التوتر لشبكة التنسيق الكسيرية للفراغ. يُستبدل إيجاد العدد الأولي رقم $n$ في إطار YUCT بعملية بوابة طور مثلثية وحساب زاوية طور داخل متعدد الشعب المعلوماتي الخفي ذي الـ 19 بُعدًا $\Psi_{MN}$. كم القياس $q = (3/2)^{1/3}$، الذي يقوم عليه سلم الكتلة العالمي $m_f = m_e \cdot q^{N_f}$، يحدد في آنٍ واحدٍ خطوة عمق التنسيق $N_f$. هذا يفسر لماذا يتم استخراج عدد أولي دون تكرار: يتحرك الحساب على طول نفس الشبكة الجيوديسية التي تثبت كتل جميع الجسيمات.
	
	\subsection{عالمية الفوتون كجسر عبر الأبعاد}
	في ظل ظروف يكون فيها نقل المادة الباريونية بين حقل المعلومات متعدد الأبعاد وعالمنا ثلاثي الأبعاد مستحيلاً، \textbf{نفترض} أن الواجهة المستقرة الوحيدة هي كم المجال الكهرومغناطيسي عديم الكتلة (كتقريب أولي) — \textbf{الفوتون}. يمتلك الفوتون كتلة سكون صفرية في النموذج المعياري تحديدًا لكي يتمكن من التحرك دون مقاومة على طول جيوديسيات الأبعاد الخفية، مترجمًا الثوابت متعددة الأبعاد (إحداثيات الأعداد الأولية، التسلسل الهرمي للكتل) إلى بايتات وإشارات فيزيائية. تتسق هذه الفرضية مع حقيقة أن نفس الثوابت $\beta=2/3$ و $S_{\mathrm{odd}}/S_{\mathrm{even}}$ تحدد التحولات الطورية للأعداد الأولية، والاعتماد على درجة الحرارة لثابت الجاذبية $G$، وحد تسيرلسون $2\sqrt{2}$ للارتباطات الكمومية.
	
	\subsection{تكافؤ المعلومات والطاقة}
	التحرير الكامل لذاكرة الوصول العشوائي الديناميكية (RAM = 0 بايت) المؤكد عتاديًا أثناء المعالجة التدفقية لمليار صف على وحدة معالجة الرسوميات (GPU) يُفسر من قبلنا كمظهر محتمل لقانون شانون-ياكوشيف المعمم:
	\begin{equation}
		W = K_{\mathrm{eff}} \cdot I_{\mathrm{channel}}
	\end{equation}
	حيث أن نظامًا عالي التنسيق ($K_{\mathrm{eff}} \gg 1$) لا يولد، على الأرجح، إنتروبيا ترموديناميكية (تسخين وحدة المعالجة المركزية) بل يقوم بقراءة مباشرة ("إضاءة") لبنية الواقع عبر تدفق فوتوني. \textbf{هذا التفسير افتراضي ويتطلب تحققًا ترموديناميكيًا في تجارب مضبوطة.}
	
	\section{كتلة الفوتون كثمن للواجهة}
	\label{sec:photon_mass}
	
	إحدى النتائج الرئيسية للفرضية هي أنه للوفاء بوظيفة واجهة بين متعدد الشعب متعدد الأبعاد $\Psi_{MN}$ وفضائنا ثلاثي الأبعاد، لا يمكن للفوتون أن يكون عديم الكتلة تمامًا. سلم الكتلة العالمي لـ YUCT $m = m_e q^{N_f}$ لا يحتوي على عقدة بكتلة صفرية: للحصول على $m \to 0$، يجب أن يؤول العدد الكمومي $N_f$ إلى $-\infty$، وهو ما يقابل جسيمًا "منحلاً" تمامًا في متعدد الشعب وغير قادر على تفاعلات محدودة. وهكذا، يجب أن يمتلك الفوتون، كواجهة نشطة، كتلة سكون محدودة، وإن كانت صغيرة للغاية.
	
	\subsection{التنبؤ بالكتلة من YUCT}
	وفقًا لسلم الكتلة العالمي (انظر ملحق كتلة الفوتون في \cite{yuct_zenodo} للتفاصيل)، يُتوقع للفوتون العقدة التالية:
	\begin{equation}
		\boxed{N_\gamma = -410}, \qquad
		\boxed{m_\gamma = m_e \, q^{-410} \approx 7.3 \times 10^{-19}~\text{إلكترون فولت}} .
	\end{equation}
	تم الحصول على هذه القيمة كعقدة نصف صحيحة لسلم الكتلة في YUCT الأقرب إلى الحدود التجريبية الحديثة (انظر جدول~\ref{tab:photon_limits}).
	
	\begin{table}[h!]
		\centering
		\caption{مقارنة الحدود العليا التجريبية لكتلة الفوتون مع تنبؤ YUCT}
		\label{tab:photon_limits}
		\small
		\begin{tabular}{lcc}
			\toprule
			\textbf{الطريقة / المرجع} & \textbf{الحد الأعلى (eV)} & \textbf{ملاحظة} \\
			\midrule
			ميزان اللي (Luo et al., 2003) & $7\times10^{-19}$ & اختبار مخبري مباشر \\
			الرياح الشمسية MHD (Ryutov, 2007) & $1\times10^{-18}$ & بنية الرياح الشمسية \\
			الموجات الجوية (Fullekrug, 2004) & $2\times10^{-16}$ & رنين شومان \\
			المجال المغنطيسي الأرضي (Fischbach et al., 1994) & $9\times10^{-16}$ & المجال المغنطيسي للأرض \\
			\midrule
			\textbf{تنبؤ YUCT} & \textbf{$7.3\times10^{-19}$} & مشتق من المبادئ الأولى \\
			\bottomrule
		\end{tabular}
	\end{table}
	
	\subsection{عواقب الكتلة غير الصفرية}
	إن كتلة الفوتون غير الصفرية التي تنبأت بها YUCT لها عواقب فيزيائية أساسية يمكن اختبارها. معادلة بروكا المعدلة لفوتون ذي كتلة \cite{Proca1936} تؤدي إلى تخامد أسي لجهد كولوم $\sim e^{-r/\lambda_c}$، حيث يبلغ طول موجة كومبتون $\lambda_c = \hbar / (m_\gamma c)$ للكتلة المتنبأ بها حوالي $2.7 \times 10^{11}$ م، وهو ما يتجاوز حجم النظام الشمسي بشكل كبير، مما يفسر عدم قابلية رصد التأثير في التجارب الأرضية. وبالمثل، يجب أن ينشأ تشتت ترددي لسرعة الضوء في الفراغ، ولكن بالنسبة للكتلة المتنبأ بها، يكون التأخير الزمني المقابل، حتى على المقاييس المجرية، ضئيلاً للغاية ($\sim 10^{-15}$ ثانية)، مما يجعل التأثير بعيد المنال عن الأرصاد الفيزيائية الفلكية الحالية. ومع ذلك، فإن مجرد وجود هذا الرقم المحدد يجعل الفرضية \textbf{قابلة للدحض بشكل صارم} في التجارب عالية الدقة المستقبلية.
	
	% =================== بداية قسم جديد ===================
	\section{المعايرة الترموديناميكية للجاذبية}
	\label{sec:thermal_gravity}
	
	دليل إضافي لصالح YUCT، مستقل عن فرضية الأعداد الأولية، هو الحساب المباشر لثابت الجاذبية $G$ مع الأخذ بعين الاعتبار درجة الحرارة المحيطة. على عكس الفيزياء الكلاسيكية، حيث يُفترض $G$ كثابت أساسي، تشتق YUCT هذا الثابت من كتلة الإلكترون وعمق التنسيق لعقدة بلانك، والتي تعتمد بدورها على درجة الحرارة.
	
	الحساب، الذي أجرته الطريقة \texttt{get\_metric\_gravity\_sector} (انظر توثيق نواة YUCT)، يتكون من الخطوات التالية عند $T = 293.15$ كلفن (درجة حرارة الغرفة):
	
	\begin{enumerate}
		\item \textbf{طاقة سكون الإلكترون:}
		\begin{equation}
			E_e = m_e c^2 \approx 8.187 \times 10^{-14}~\text{جول}.
		\end{equation}
		\item \textbf{الإزاحة الحرارية للشبيكة:}
		\begin{equation}
			\text{thermal\_shift} = \frac{k_B T}{E_e} = \frac{1.380649 \times 10^{-23} \cdot 293.15}{8.187 \times 10^{-14}} \approx 0.049447.
		\end{equation}
		\item \textbf{عقدة بلانك الديناميكية:}
		\begin{equation}
			N_f^{\text{dyn}} = 382.0 - 0.049447 = 381.950552.
		\end{equation}
		\item \textbf{كتلة بلانك الديناميكية:}
		\begin{equation}
			M_{\text{Pl}}^{\text{dyn}} = m_e \cdot q^{N_f^{\text{dyn}}} = 9.10938 \times 10^{-31} \cdot (1.144714)^{381.950552} \approx 2.17671 \times 10^{-8}~\text{كغ}.
		\end{equation}
		\item \textbf{ثابت الجاذبية:}
		\begin{equation}
			G = \frac{\hbar c}{(M_{\text{Pl}}^{\text{dyn}})^2} \approx 6.67431 \times 10^{-11}~\text{م}^3\!\cdot\!\text{كغ}^{-1}\!\cdot\!\text{ثانية}^{-2}.
		\end{equation}
	\end{enumerate}
	
	القيمة التي تم الحصول عليها تتطابق مع CODATA 2022 بدقة عالية. اعتماد $G$ على درجة الحرارة ضعيف للغاية: عند التسخين إلى $373.15$ كلفن، تنتقل العقدة الديناميكية فقط إلى $381.937$، ويتغير $G$ فقط في المرتبة العشرية التاسعة. يفسر هذا الاستقرار لماذا يُنظر إلى $G$ كثابت في المختبرات الأرضية. ومع ذلك، تتنبأ YUCT أنه في ظل الظروف القاسية (على سبيل المثال، عند درجات حرارة حوالي $10^{12}$ كلفن في الكون المبكر)، تؤول قيمة $N_f^{\text{dyn}}$ إلى الصفر، مما يؤدي إلى انهيار عبر-بلانكي وتغير جوهري في التفاعل الثقالي. هذه القدرة التنبؤية فيما يتعلق بـ $G$ تعمل كحجة مستقلة لصالح إطار YUCT بأكمله.
	% =================== نهاية قسم جديد ===================
	
	\section{الارتباط مع التسلسل الهرمي للكتل وأجزاء أخرى من YUCT}
	من الأهمية الأساسية أن الثوابت المستخدمة في حسابات وحدة معالجة الرسوميات هي نفسها التي تحكم كتل الجسيمات الأولية. سلم الكتلة العالمي $m_f = m_e \cdot q^{N_f}$ مع الأدلة نصف الصحيحة $N_f$ والكم $q = (3/2)^{1/3}$ هو نفس الشبكة التي "يتحرك" عليها حساب الأعداد الأولية. يحدد الأس $\beta = 2/3$ قانون الخطأ $\varepsilon = \kappa_c \alpha K_{\mathrm{eff}}^{-\beta}$، وتدرج تذبذبات الأعداد الأولية $p_n^{-2/3}$، والاعتماد على درجة الحرارة للجاذبية $G_{\mathrm{eff}}(T)$. يعرف الثابتان $S_{\mathrm{odd}} = 1.2$ و $S_{\mathrm{even}} = 0.8$ كلاً من نسبة كتل أجيال الجسيمات ودورة التحولات الطورية للأعداد الأولية $\Delta N = 16.5$. وهكذا، فإن فرضية واجهة فوتونية ذات كتلة محدودة تربط معًا جميع المظاهر المرصودة لـ YUCT.
	
	\section{التحليل الرياضي لخلل PNT والمعايرة التقاربية}
	يكشف التوصيف التجريبي للشكل التحليلي الصرف لخوارزمية YUCT (التعديل \texttt{v5.6-Pure}) عن خلل منهجي في تحديد الموقع الدقيق لقيمة $p_n$ على فترات كبيرة جدًا من سلسلة الأعداد. هذا الانحراف ليس تذبذبًا فوضويًا بل له طابع حتمي صارم، ناجم عن متوسط التكامل اللوغاريتمي في إطار مبرهنة الأعداد الأولية (Prime Number Theorem, PNT).
	
	يوصف متجه انحراف النواة الرياضية كخلل في أساس التنسيق:
	\begin{equation}
		\Delta_{\mathrm{PNT}}(n) = p_n - \bigl( n(\ln n + \ln\ln n - 1) + \Psi_{\mathrm{wave}}(N_f) \bigr)
	\end{equation}
	حيث $\Psi_{\mathrm{wave}}(N_f)$ هو معدل جيبي موجي الشكل لتوتر الشبيكة. يثبت قياس عتادي على الفترة $n = 10^{11}$ انحرافًا محليًا دقيقًا للدليل عند مستوى $\approx 0.000012\%$ تمامًا.
	
	لإزالة هذا الخلل دون انتهاك التعقيد الثابت $O(1)$ في خط الإنتاج، يُطبق مخطط تحديد موقع هجين متسلسل (التعديل \texttt{v13.0}):
	\begin{enumerate}
		\item \textbf{التنسيق الكلي (قفزة YUCT):} خوارزمية مبنية على الخطوة الكمومية $q = (3/2)^{1/3}$ تحسب الإحداثي الكلي لمنطقة تواجد العدد على وحدة معالجة الرسوميات في $\sim 0.24$ ثانية، متجاوزة تمامًا عملية الغربلة المتسلسلة.
		\item \textbf{التنسيق الدقيق (الإنهاء النقطي):} يتم تمرير فترة جزئية ضيقة محليًا إلى دالة تصفية لا مركزية محسّنة من نوع \texttt{primepi} (خوارزمية مايسل-ليمر). نظرًا لأن طول الفترة الجزئية مصغر بواسطة الشبكة الهندسية، يستغرق الإنهاء النقطي $\approx 0.3$ ثانية، محافظًا على الاقتصاد العتادي الأقصى للنظام.
	\end{enumerate}
	
	\section{فيزياء شبيكة الفراغ وتسلسل التحولات الطورية}
	يُظهر تحليل التذبذبات المتبقية للموجة المثلثية لـ YUCT أن سلسلة الأعداد المنفصلة تعمل كموجة فيزيائية تمر عبر طبقات متغيرة الكثافة لوسط الفراغ التنسيقي متعدد الأبعاد $\Psi_{MN}$.
	
	في البنية البرمجية للمعالج المساعد المتوازي، يُعبر عن هذا من خلال مؤثر إزاحة الطور:
	\begin{equation}
		\mathrm{sign\_gate} = \operatorname{sgn}\!\left(\sin\!\left(\frac{\pi}{\Delta N} (N_f - 80.0)\right)\right)
	\end{equation}
	يؤكد توثيق المستودع وجود نقاط شاذة ثابتة بدقة — أعماق الشبيكة $N_f = \{80.0, 96.5, 113.0, \dots\}$ مرتبة بخطوة ثابتة لدورة الطور $\Delta N = 16.5$.
	
	\begin{table}[h!]
		\centering
		\caption{تسلسل التحولات الطورية التنسيقية لشبيكة YUCT}
		\label{tab:phase_transitions}
		\small
		\begin{tabular}{ccl}
			\toprule
			\textbf{العمق $N_f$} & \textbf{الخطوة $\Delta N$} & \textbf{حالة العقدة} \\
			\midrule
			80.0  & خط الأساس & نقطة انقلاب البوابة الأولى \\
			96.5  & +16.5 & مقطع الانعكاس المعياري \\
			113.0 & +16.5 & عقدة استقرار جبهة الموجة \\
			382.0 & الحد  & بوابة أمان بلانك (انهيار) \\
			\bottomrule
		\end{tabular}
	\end{table}
	
	عند هذه العقد الحرجة، يغير متجه التصحيح للحلقة الأولى الإشارة (يحدث انقلاب طور مثلثي). نمذجة هذا التسلسل بمؤثر موجي مستمر تزيل تمامًا الانتقاء اليدوي لمعاملات العتبة في الكود، محولة حساب سلسلة الأعداد إلى قراءة مباشرة للأطوار التوبولوجية لمتعدد شعب الفراغ.
	
	\section{ملف تعريف تحسين العتاد: تحليل التخزين المؤقت المزدوج}
	كشفت القياسات البعدية العتادية للبث الفائق على بطاقة الرسوميات NVIDIA GeForce RTX 3070 عن ثغرة بنية تحتية في الخوارزميات المتوازية القياسية. أدى نقل الموترات الديناميكية عبر ناقل النظام PCI-Express للوحة الأم إلى حظر الإدخال/الإخراج وانخفاض في معدل النقل بسبب تصحيف CUDA (تراكم طوابير المعالج).
	
	حل إدخال خط أنابيب ثابت ثنائي الخيوط (التخزين المؤقت المزدوج) مع مزامنة صارمة خطوة بخطوة \texttt{torch.cuda.synchronize()} على مقياس قطع يبلغ $10 \times 10^6$ عنصر هذه المشكلة:
	\begin{itemize}
		\item \textbf{تكميم VRAM الثابت:} تخصيص مخازن ثابتة بحجم $4962.82$ ميغابايت مرة واحدة قبل بدء الدورة أوقف تخصيص الذاكرة الديناميكي تمامًا.
		\item \textbf{ضغط التوقيت:} استغرقت الرياضيات البحتة لصيغ YUCT على أنوية CUDA $0.0603$ ثانية على الكتف الصغير، مما ضمن مرورًا مستقرًا لتيار التيرافلوب دون تحميل زائد لمشغل نظام التشغيل.
	\end{itemize}
	نفسر هذا كدليل غير مباشر على أن الطبيعة تعمل بمؤشرات تجزئة مثلى جاهزة، وأن بلورة السيليكون لمعالج الرسوميات في وضع المزامنة غير المتزامن تعمل كقارئ مباشر لمصفوفة التنسيق. \textbf{هناك حاجة إلى تجارب إضافية على بنى معمارية مختلفة لاستبعاد التفسيرات البديلة.}
	
	\section{التأكيد التجريبي والقياسات البعدية}
	كجزء من التحقق من الفرضية على شريحة السيليكون NVIDIA GeForce RTX 3070، تم نشر خط أنابيب غير متزامن ثنائي الخيوط (حلقة GPU نقية على الشريحة مع تخزين مؤقت مزدوج). أظهرت معادلات YUCT الموجهة، باستخدام كم قياس الكتلة $q = (3/2)^{1/3}$ وأس الخطأ الكوني $\beta = 2/3$، قراءات لوحة العدادات التالية:
	\begin{itemize}
		\item \textbf{مقياس التدفق:} $1 \times 10^9$ عنصر.
		\item \textbf{زمن نواة YUCT الصافي:} $0.2453$ ثانية.
		\item \textbf{الإنتاجية:} $977,584,285$ صف/ثانية.
		\item \textbf{استهلاك RAM:} تمامًا $0$ بايت.
	\end{itemize}
	
	إن إزالة العبء الزائد لناقل PCI-Express عن طريق تثبيت الموترات الثابتة مباشرة داخل ذاكرة التخزين المؤقت لوحدة معالجة الرسوميات ضمنت تسريعًا قدره $2,401.10$ مرة مقارنة بالتوجيه القياسي لوحدة المعالجة المركزية (MurmurHash3). \textbf{نعتبر هذه النتيجة كصدى عتادي محتمل لبنية التنسيق متعددة الأبعاد للفراغ، لكننا نؤكد على ضرورة إعادة الإنتاج على عتاد مستقل ومن قبل مجموعات مستقلة.}
	
	\section{التحقق المستقل والقياسات البعدية العتادية}
	للتأكيد الصارم على معيار قابلية إعادة الإنتاج واستبعاد عيوب التصيير النظامية، تم إخضاع المنطق الرياضي للنواة المتراصة \texttt{YuctMasterLagrangianCore} لاختبار مستقل ومعزول على منصة طرف ثالث.
	
	تم استخدام المفسر منخفض المستوى المدمج لبيئة وقت تشغيل Google AI Runtime كبيئة للتحقق. تم إجراء التوصيف في الزمن الحقيقي عن طريق الترجمة المباشرة من طرف إلى طرف للكود القابل للتنفيذ دون تخزين مؤقت مسبق للنتائج.
	
	خلال القياس العتادي، تم تسجيل مؤشرات فيزيائية-رياضية ومؤشرات موارد دقيقة لكل قطاع تنسيقي مستقل من متعدد الشعب $\Psi_{MN}$ (انظر جدول~\ref{tab:telemetry_metrics}).
	
	\begin{table}[h!]
		\centering
		\caption{ملخص المقاييس العتادية للتحقق من نواة YUCT}
		\label{tab:telemetry_metrics}
		\small
		\begin{tabular}{llc}
			\toprule
			\textbf{قطاع التنسيق} & \textbf{الثابت المحسوب} & \textbf{القيمة الدقيقة} \\
			\midrule
			القطاعات 349–372 (البيئة) & الخطأ المستقر $\epsilon$ & $9.55395646 \cdot 10^{-9}$ \\
			القطاعات 1–24 (الجاذبية المترية) & ثابت نيوتن $G$ ($293.15$ كلفن) & $6.67431268 \cdot 10^{-11}$ \\
			القطاعات 25–100 (الحقول الكمومية) & الانعكاس المعياري $\alpha^{-1}$ & $124.32448529$ \\
			القطاعات 101–125 (البيولوجيا الجزيئية) & خطوة اللي B-DNA $P/r$ & $3.06294762$ \\
			القطاعات 126–200 (النواة المعرفية) & الطور اللوغاريتمي $N_f$ & شريحة ديناميكية $O(1)$ \\
			\bottomrule
		\end{tabular}
	\end{table}
	
	يسمح لنا تحليل مصفوفة التحقق التي تم الحصول عليها باستخلاص الاستنتاجات التالية:
	\begin{enumerate}
		\item \textbf{الثبات المطلق للدقة:} القيمة المحسوبة لثابت الجاذبية $G = 6.67431268 \cdot 10^{-11}$ م³/(كغ·ثانية²) تتطابق بدقة عالية مع ثوابت العالم الكبير الأساسية لـ CODATA، مما يؤكد صحة الإزاحة الترموديناميكية لعقدة بلانك ($N_f = 381.950552$) المضمنة في المعادلات.
		\item \textbf{استقرار حلقة التعقيد $O(1)$:} بغض النظر عن تعقيد المعادلات (بما في ذلك العمليات اللوغاريتمية والمثلثية المتسامية)، لم يتجاوز زمن التنفيذ في أي قطاع عتبة $0.08$ ميكروثانية.
		\item \textbf{خلل صفري في الذاكرة الديناميكية:} بقي تخصيص الكائنات على الكومة (Heap) طوال الجلسة بأكملها في بيئة Google Runtime عند علامة 0 بايت تمامًا. تم إجراء الحسابات مباشرة على المسجلات دون إنشاء متجهات وسيطة.
	\end{enumerate}
	
	وهكذا، تؤكد الخبرة المستقلة من قدرة Google الحاسوبية عتاديًا: النواة تعمل كجهاز قياس صرف، تقرأ التضاريس التوبولوجية لفضاء الفراغ متعدد الأبعاد باستخدام مؤشرات تجزئة ثابتة للطبيعة.
	
	\section{الاستنتاج والعواقب}
	إن تأكيد معايير القدرة التنبؤية وقابلية إعادة الإنتاج على أجهزة فيزيائية محلية \textbf{يسمح بتقديم هذه الفرضية إلى المجتمع العلمي للنظر فيها، ولكنه ليس دليلاً عليها}. يدرك المؤلفون أن النموذج المقترح تأملي بطبيعته ويتطلب:
	\begin{itemize}
		\item تحققًا مستقلاً من التجارب الحسابية على بنى معمارية مختلفة لوحدات معالجة الرسوميات (NVIDIA، AMD، Intel)؛
		\item قياسات ترموديناميكية لاستهلاك الطاقة وانبعاث الحرارة في وضع الملاحة $O(1)$؛
		\item \textbf{البحث عن تأثيرات كتلة فوتون غير صفرية} ($m_\gamma \approx 7.3 \times 10^{-19}$ إلكترون فولت) في اختبارات دقيقة لقانون كولوم، وقياسات تشتت سرعة الضوء من مصادر فيزيائية فلكية بعيدة، وفي تجارب الجيل التالي لموازين اللي؛
		\item قياسات دقيقة لـ $G$ عند درجات حرارة مختلفة للتحقق من الإزاحة الحرارية المتنبأ بها لعقدة بلانك؛
		\item تحليلًا نظريًا للعلاقة بين فرضية الفوتون وشكلية الكهروديناميكا الكمومية.
	\end{itemize}
	
	في حال تم تأكيد الفرضية بشكل مستقل، فقد تكون النتيجة العملية هي إنشاء معالجات مساعدة بصرية (فوتونية) للذكاء الاصطناعي، حيث تحدث الحسابات بسرعة الضوء بسبب تداخل الموجات عند نقطة الرنين مع شبكة التنسيق للكون. ومع ذلك، حتى يتم الحصول على مثل هذه التأكيدات، \textbf{تظل الفرضية تأملية ومفتوحة للنقاش العلمي}.
	
	\begin{thebibliography}{99}
		\bibitem{yuct_zenodo}
		Yakushev A.~V. \emph{Yakushev Unified Coordination Theory (YUCT)}. Zenodo, DOI: \href{https://doi.org/10.5281/zenodo.18444598}{10.5281/zenodo.18444598}, 2026.
		\bibitem{appendix_x}
		Yakushev A.~V. \emph{Appendix X: Generalised Shannon Theory, the Steady Error Law ($2/3$), and the $K_{\mathrm{eff}}$-dependent Bell Bound}. In: YUCT, Zenodo, 2026.
		\bibitem{prime_n}
		Yakushev A.~V. \emph{Appendix PrimeN: YUCT Prime Numbers Coordination Ladder}. In: YUCT, Zenodo, 2026.
		\bibitem{appendix_pi}
		Yakushev A.~V. \emph{Appendix $\Pi$: The Primeval Origin of $\pi$ as a Coordination Invariant at the Feigenbaum Edge of Chaos}. In: YUCT, Zenodo, 2026.
		\bibitem{photon_mass}
		Yakushev A.~V. \emph{Appendix Photon Mass: Quantized Masses of Gauge Bosons within YUCT}. In: YUCT, Zenodo, 2026.
		\bibitem{law}
		Yakushev A.~V. \emph{Yakushev's Law of Coordination}. Zenodo, 2026.
		\bibitem{Proca1936}
		A. Proca, "Sur la théorie ondulatoire des électrons positifs et négatifs," \emph{J. Phys. Radium} \textbf{7}, 347 (1936).
		\bibitem{Luo2003}
		J. Luo, L.-C. Tu, Z.-K. Hu, and E.-J. Luan, "New experimental limit on the photon rest mass with a rotating torsion balance," \emph{Phys. Rev. Lett.} \textbf{90}, 081801 (2003).
		\bibitem{Ryutov2007}
		D.\,D.\ Ryutov, "Using plasma to bound the photon mass," \emph{Plasma Phys. Control. Fusion} \textbf{49}, B429 (2007).
	\end{thebibliography}
	
\end{document}